%% ---------------------------------------------------------------------------
%% Front matter, introduction, related work, problem formulation, methods.
%% Results-dependent sections are in results.tex.
%% ---------------------------------------------------------------------------

\section{Introduction}
\label{sec:intro}

Most applied machine learning on business, health and public-sector records still works with tabular data, and its modelling toolkit is mature. Gradient-boosted trees remain the strongest general baseline on medium-sized tables~\cite{grinsztajn2022}, living benchmarks track dozens of competing learners~\cite{erickson2025}, and pre-trained tabular foundation models now make strong predictions on small datasets in a single forward pass~\cite{hollmann2025}. What the toolkit does not settle is a question that comes before modelling: \emph{is there anything in this table to learn?} When the target is nearly independent of the recorded features, no amount of tuning can rescue the project, yet teams typically discover this only after weeks of cleaning, feature engineering and tuning.

Work on automatic feasibility studies has framed this as a separate step that estimates, before development starts, whether a target accuracy is reachable~\cite{renggli2020snoopy,renggli2023feasibility}. That work estimates the \emph{size} of the best achievable error and concentrates on image and text data. The narrower question for small and medium tabular data---whether a dataset carries \emph{any} detectable learnable signal, and how far a negative answer can be trusted---has not been studied systematically. Many inexpensive probes could answer it: information-theoretic and statistical dependence measures, geometric data-complexity measures, and cheap baseline learners. They come from separate traditions and are reported on their own terms, without a common definition of signal, a common inference procedure, or a known false-alarm rate~\cite{lorena2019,ross2014,ojala2010}.

The obvious check---train a model and compare it with a majority-class baseline---is less informative than it looks. Cross-validated scores on a few hundred rows carry error bars of about $\pm$10 percentage points~\cite{varoquaux2018}, and a single failed model cannot distinguish ``no signal'' from ``wrong model class''. This paper therefore compares eleven probes from three families under one definition of signal, one permutation-based inference layer and one ground truth. Its contributions are:
\begin{enumerate}
\item a shared inference layer that turns statistical, geometric and model-based probes into calibrated tests, with Benjamini--Hochberg/Simes aggregation for per-feature tests~\cite{korthauer2019};
\item a controlled benchmark of false-positive rate, power and minimum detectable effect across sample sizes, signal types, irrelevant features and class imbalance, using synthetic data with known Bayes AUC, label-shuffled real data, and two randomised experiments intended to be null by design;
\item a replication of the two permutation tests of Ojala and Garriga~\cite{ojala2010} on every dataset, which separates signal located in single features from signal located in feature dependencies;
\item a two-stage go/no-go protocol, designed on simulated data and validated on held-out data, and an open-source implementation in documented Jupyter notebooks.
\end{enumerate}

\section{Related Work}
\label{sec:related}

\textbf{Feasibility studies and Bayes-error estimation.} Ease.ml/snoopy estimates the Bayes error from cheap proxies to decide whether a user-specified accuracy is attainable~\cite{renggli2020snoopy}, and follow-up work relates that estimate to label noise~\cite{renggli2023feasibility}. FeeBee evaluates Bayes-error estimators by injecting known label noise into real data~\cite{renggli2021feebee}; Theisen et al.\ compute exact Bayes errors for data generated by normalising flows~\cite{theisen2021}; Ishida et al.\ estimate it directly from soft labels~\cite{ishida2023}. These studies estimate how good a model \emph{could} be. Detecting whether any signal exists is a weaker question that admits exact error control, which is the angle taken here.

\textbf{Data-complexity measures.} Lorena et al.\ standardised more than twenty measures of classification difficulty, including the maximum Fisher's discriminant ratio (F1), the fraction of borderline points on a minimum spanning tree (N1) and the leave-one-out 1-NN error (N3)~\cite{lorena2019}. Instance hardness~\cite{smith2014}, instance-space analysis~\cite{munoz2018} and meta-learning~\cite{rivolli2022} build on them, and the \texttt{problexity} library implements them in Python~\cite{komorniczak2022}. These measures are descriptive: they carry no null distribution, so a value such as N1 = 0.4 has no fixed meaning across sample sizes and dimensions.

\textbf{Dependence measures.} Mutual information (MI) captures any dependence, but its estimation is hard. scikit-learn~\cite{pedregosa2011} uses the nearest-neighbour estimator of Ross~\cite{ross2014}; Gao et al.\ showed that such estimators can fail on discrete--continuous mixtures~\cite{gao2017}; McAllester and Stratos proved that high-confidence lower bounds on MI cannot exceed about $\log n$~\cite{mcallester2020}; and Czy\.z et al.\ found estimator accuracy to depend strongly on the data distribution~\cite{czyz2023}. Distance correlation~\cite{szekely2013,ramoscarreno2023} and the Hilbert--Schmidt independence criterion (HSIC) test independence between a whole vector and the label; they belong to one family of statistics~\cite{sejdinovic2013}, and their power can decay with dimension~\cite{ramdas2015}. The dependence literature mostly serves feature selection~\cite{vergara2014,li2017}, where combining many per-feature tests without correction inflates false positives~\cite{korthauer2019}.

\textbf{Classifier-based tests.} For a binary label, testing $Y\perp\!\!\!\perp X$ is equivalent to testing whether the two class-conditional distributions are equal, which links signal detection to classifier two-sample tests~\cite{lopezpaz2017,kim2021}. Ojala and Garriga formalised permutation tests for cross-validated classifier error~\cite{ojala2010}; permutation $p$-values with randomly drawn permutations are exact when computed with the ``$+1$'' correction~\cite{hemerik2018,phipson2010}.

\textbf{Foundation models, evaluation and leakage.} TabPFN~\cite{hollmann2023,hollmann2025} and TabICL~\cite{qu2025} perform in-context prediction on tables, and language models can classify serialised rows~\cite{hegselmann2023}, although they have memorised popular tabular datasets~\cite{bordt2024}. Cross-validation estimates average performance over training sets rather than the fitted model's error~\cite{bates2024}, benchmark results vary with seeds and splits~\cite{bouthillier2021}, learning curves inform data-collection decisions~\cite{mohr2024,viering2023}, and leakage is widespread in published machine-learning-based science~\cite{kaufman2012,kapoor2023}. Label-noise research~\cite{frenay2014} informs the label-shuffling design used below.

\section{Problem Formulation}
\label{sec:problem}

Let $D=\{(x_i,y_i)\}_{i=1}^n$ be $n$ exchangeable draws of a feature vector $X$ (numeric and categorical) and a binary label $Y$. A dataset has \emph{no signal} when
\begin{equation}
H_0:\; Y\perp\!\!\!\perp X \;\Longleftrightarrow\; P(Y\mid X)=P(Y) \;\Longleftrightarrow\; I(X;Y)=0 .
\label{eq:h0}
\end{equation}
For a loss $\ell$ with risk $R(f)$, the population signal is the Bayes gap
\begin{equation}
\Delta^{*}=R(f_{\mathrm{const}})-R(f^{*})\ge 0,
\label{eq:bayesgap}
\end{equation}
where $f_{\mathrm{const}}$ always predicts the majority class and $f^{*}$ is Bayes-optimal; $\Delta^{*}=0$ under $H_0$. Under 0--1 loss the converse can fail (a weak dependence that never changes the most likely class leaves $\Delta^*=0$), which is one reason this study measures signal with the area under the ROC curve (AUC). Because only a restricted model family $\mathcal{V}$ can be trained on $n$ samples, the practically relevant quantity is closer to \emph{predictive $\mathcal{V}$-information}~\cite{xu2020,ethayarajh2022}, the reduction in achievable log loss when models in $\mathcal{V}$ may use $X$. A model-based probe estimates such a gap for its own family; MI is the limit in which $\mathcal{V}$ contains every model.

Two consequences shape the evaluation. First, \emph{significance and usefulness differ}: with $n=50{,}000$ an AUC of 0.51 can be significant yet useless, so the protocol reports an effect size with its interval beside every $p$-value. Second, \emph{verdicts are asymmetric}: rejecting $H_0$ is evidence of signal, but a non-rejection may reflect low power or a probe that is blind to the structure present. The study therefore reports ``no signal detected'' together with a minimum detectable effect (MDE), the smallest Bayes AUC at which the probe reaches 80\% power.

\section{Methods}
\label{sec:methods}

\subsection{Probes}

Table~\ref{tab:probes} lists the eleven probes. Each returns a statistic oriented so that larger values mean stronger evidence of signal. All preprocessing is unsupervised (median/mode imputation, standardisation, one-hot encoding), so it is identical for the observed and every permuted dataset and leaks no label information into cross-validation.

\begin{table}[!t]
\centering
\caption{Probes compared in this study.}
\label{tab:probes}
\footnotesize
\begin{tabular}{p{0.9cm}p{2.6cm}p{3.9cm}}
\hline
ID & Probe & Statistic (larger = more signal) \\ \hline
S1 & Mutual information & per-feature MI (Ross estimator~\cite{ross2014}); Simes-combined \\
S2 & ANOVA F / $\chi^2$ & per-feature analytic $p$-values; Simes-combined \\
S3 & Distance correlation & $\mathrm{dCor}(X,Y)$~\cite{szekely2013} \\
S4 & HSIC & Gaussian kernel, median bandwidth \\
G1 & Fisher ratio F1 & $1-\mathrm{F1}$ (\texttt{problexity}~\cite{komorniczak2022}) \\
G2 & Borderline points N1 & $-\mathrm{N1}$ on the Euclidean MST \\
G3 & 1-NN error N3 & $-\mathrm{N3}$ (leave-one-out) \\
M1 & Decision tree & CV AUC, depth $\le$ 3 \\
M2 & Logistic regression & CV AUC, L2, $C=1$ \\
M3 & Random forest & CV AUC, 50 trees, leaf size 5 \\
M4 & Gradient boosting & CV AUC, histogram GBDT (reference) \\ \hline
\end{tabular}
\par\vspace{2pt}\parbox{0.95\linewidth}{\scriptsize CV AUC: out-of-fold AUC from stratified 5-fold cross-validation; the AUC gain over the dummy baseline (AUC 0.5) is the effect size.}
\end{table}

For a binary label the two kernel-type probes reduce to quadratic forms. With $A$ the double-centred Euclidean distance matrix of $X$, whose rows sum to zero,
\begin{equation}
\mathrm{dCov}^2_n(X,Y)=\frac{1}{n^2}\sum_{ij}A_{ij}\,|y_i-y_j|=-\frac{2}{n^2}\,y^{\top}Ay ,
\end{equation}
and with $\tilde K$ the centred Gaussian kernel matrix and the delta kernel on $Y$, $\mathrm{HSIC}_b=\tfrac{2}{n^2}\,y^{\top}\tilde K y$. Pre-computing $A$, $\tilde K$, the minimum spanning tree (G2) and the nearest-neighbour graph (G3) once per dataset makes each permutation a single matrix--vector product or a vector comparison. The fast implementations were verified against textbook formulas and against \texttt{problexity}.

\subsection{Shared Inference Layer}

Every probe is wrapped in a label-permutation test~\cite{ojala2010}: labels are permuted $B$ times, the full probe including cross-validation is recomputed, and
\begin{equation}
p=\frac{1+\#\{b: T_b\ge T_{\mathrm{obs}}\}}{B+1},
\end{equation}
which is exact under exchangeability~\cite{hemerik2018,phipson2010}. For per-feature probes, the dataset-level $p$-value is the Simes combination $\min_k m\,p_{(k)}/k$; rejecting when it is at most $\alpha$ is identical to the Benjamini--Hochberg procedure rejecting at least one feature~\cite{korthauer2019}. The uncorrected rule (``any $p_j\le\alpha$'') is recorded to measure its inflation. In synthetic experiments, where the generator is known, the null distribution of each statistic is simulated once per setting from reference datasets with $Y$ independent of $X$ (a Monte Carlo null), and held-out null datasets measure the resulting false-positive rate. $\alpha=0.05$ throughout.

\subsection{Ojala--Garriga Tests}

Test~1 permutes the labels and tests $H_0$ of~\eqref{eq:h0}. Test~2 permutes each feature independently \emph{within each class}, which preserves every class-conditional marginal but destroys dependencies between features; a small $p$-value indicates that the classifier exploits feature dependencies~\cite{ojala2010}. Following the original paper, the error is estimated by 10-fold cross-validation on fixed folds, and $p=(\#\{D': e(f,D')\le e(f,D)\}+1)/(k+1)$. Decision tree, Gaussian naive Bayes, 1-NN and RBF-SVM classifiers are used with $k=200$ permutations per test. A categorical feature's one-hot columns are permuted as one block.

\subsection{Two-Stage Protocol}

The protocol first runs a sanity check for identifier-like columns, exact duplicates and any single feature with AUC above 0.99 (``leakage suspected''). Stage~1 is the better of the two multivariate dependence tests (S3, S4) and stage~2 the best model-based probe, both chosen by mean power in Experiment~B only. $H_0$ is rejected if either stage gives $p\le\alpha/2$ (Bonferroni), and every verdict reports the AUC gain with its interval or, for a non-rejection, the MDE.
